%% kapa: exam=1 round=30 year=2019 period=2 form=A subject=appraisal_law subject_ko=감정평가관계법규 seq=9 number=9
%% source: Q-Net 감정평가사 기출문제 (cst00309) artlSeq=5205205 file=extracted/1008/5205205/2019년 제30회 감정평가사 1차시험 문제지/2019년 제30회 감정평가사 1차시험 문제지/2019년도 제30회 감정평가사 1차 2교시 A형.hwp
%% conversion: kordoc 4.24.51 -> tools/gen1.py (auto); review=proofread-sonnet; flags=none
\begin{question}{9}
국토의 계획 및 이용에 관한 법령상 용도지역에 관한 설명이다. ( )에 들어갈 용어가 옳게 연결된 것은?
\begin{framedtext}
○ ( ㄱ ): 중고층주택을 중심으로 편리한 주거환경을 조성하기 위하여 필요한 지역 \par
○ ( ㄴ ): 환경을 저해하지 아니하는 공업의 배치를 위하여 필요한 지역 \par
○ ( ㄷ ): 도시의 녹지공간의 확보, 도시확산의 방지, 장래 도시용지의 공급 등을 위하여 보전할 필요가 있는 지역으로서 불가피한 경우에 한하여 제한적인 개발이 허용되는 지역 \par
\end{framedtext}
\begin{choices}
  \item ㄱ: 제2종일반주거지역, ㄴ: 준공업지역, ㄷ: 자연녹지지역
  \item ㄱ: 제2종일반주거지역, ㄴ: 준공업지역, ㄷ: 보전녹지지역
  \item ㄱ: 제2종일반주거지역, ㄴ: 일반공업지역, ㄷ: 자연녹지지역
  \item ㄱ: 제3종일반주거지역, ㄴ: 일반공업지역, ㄷ: 보전녹지지역
  \item ㄱ: 제3종일반주거지역, ㄴ: 일반공업지역, ㄷ: 자연녹지지역
\end{choices}
\end{question}
\answer{5}
